\section{引言}

在大规模图像识别中，网络深度是精度的重要来源：VGG 使用 3x3 小卷积核，把深度推到 16-19 个权重层后显著优于此前的配置~\cite{simonyan2015}。然而，深度增加后训练变得更加困难，非常深的网络如何训练仍是开放问题~\cite{srivastava2015}。因此，能否稳定地训练比现有网络深得多的模型，是继续利用深度收益的前提。

\section{相关工作}

\subsection{通过加深网络提升表示能力}

这一路线的共同做法是在同一识别任务上不断增加层数，区别在于增加深度的方式~\cite{simonyan2015,szegedy2015}。VGG 保持结构简单，只堆叠小卷积核~\cite{simonyan2015}；GoogLeNet 则采用更深的模块化结构~\cite{szegedy2015}。二者都没有改变“逐层堆叠”的训练方式。

\subsection{改善深层网络的可训练性}

另一条路线不改变深度，而是改变训练信号穿过多层的方式~\cite{ioffe2015,srivastava2015}。批归一化对层输入做归一化，使网络可以使用高得多的学习率并降低对初始化的依赖~\cite{ioffe2015}，在图像分类模型上以少 14 倍的训练步数达到相同精度~\cite{ioffe2015}。与之不同，Highway 网络借鉴 LSTM，用自适应门控单元调节跨层信息流~\cite{srivastava2015}。前者作用于激活分布，后者引入了带参数的捷径。

上述方法解决了深层网络能否开始收敛的问题~\cite{ioffe2015}，但门控捷径依赖数据相关的参数~\cite{srivastava2015}，更深网络的训练本身仍是开放问题~\cite{srivastava2015}。尚缺少一种不增加参数、又能让更深网络至少不比浅层网络更难优化的结构。

本文的贡献如下：

\begin{itemize}
\item 提出残差学习框架，用不增加参数和计算量的恒等捷径连接让堆叠层拟合残差映射 (图2)。
\item 在 ImageNet 上训练深达 152 层的残差网络，深度为 VGG 的 8 倍但复杂度更低，集成模型在测试集上的 top-5 错误率为 3.57\% (表5)。
\item 在 CIFAR-10 上分析 100 层与 1000 层以上的残差网络，并与 FitNet、Highway 比较 (表6)。
\end{itemize}
